\documentclass[UTF8]{ctexrep}
\usepackage{amsmath,amssymb,bm}
\title{科学公式结构回归}
\author{Serein}
\date{}

\begin{document}
\maketitle

% 该注释不能出现在导出正文。
\chapter{基础结构}

\section{分段函数与矩阵}

\noindent 分段函数和矩阵应保持为原生数学结构：

\[
f(x)
=
\begin{cases}
x^2, & x\geq 0,\\
-x, & x<0,
\end{cases}
\qquad
\bm A
=
\begin{bmatrix}
a_{11} & a_{12}\\
a_{21} & a_{22}
\end{bmatrix}.
\]

\subsection{长公式与集中公式}

\begin{multline}
J(\bm\theta)
=
\sum_{i=1}^{n}
\left[y_i-f(\bm x_i;\bm\theta)\right]^2\\
+\lambda\sum_{j=1}^{p}\theta_j^2.
\label{eq:structures-multline}
\end{multline}

\begin{gather*}
a^2+b^2=c^2,\\
e^{i\pi}+1=0,\\
\sin^2x+\cos^2x=1.
\end{gather*}

\chapter{编号与引用}

\section{子公式}

\begin{subequations}
\label{eq:structures-system}
\begin{align}
\bm g_k&=\bm J_k^{\mathsf T}\bm r_k,
\label{eq:structures-gradient}\\
\bm H_k&\approx\bm J_k^{\mathsf T}\bm J_k.
\label{eq:structures-hessian}
\end{align}
\end{subequations}

引用式~\eqref{eq:structures-gradient} 和式~\eqref{eq:structures-hessian}。

\subsection{手动编号}

\begin{equation}
E=mc^2
\tag{A.1}
\end{equation}

\subsubsection{无编号公式}

\begin{equation*}
\operatorname{diag}(\bm A)
=
\begin{bmatrix}a_{11}&a_{22}\end{bmatrix}.
\end{equation*}

\paragraph{结论}

上述样本覆盖标题层级、矩阵、分段函数、多行公式、子公式和引用。

\end{document}
